% !PRINTED 590
\weberpage{590}
%<<CONT
the $\varkappa$, apart from unit factors, must agree with the totality
of the $\pi$.

It now remains to actually determine the primes of the
field $R(i)$.

Every prime is a divisor of infinitely many rational
integers, e.g. of the norm and all its multiples.
Among the rational positive integers that are divisible by a prime
$\pi$ there is one, which we shall denote by $n$,
that is the smallest, and this must be a prime in the field $R$.
For if it were decomposable in $R$, then by 5. one of its factors,
and hence a still smaller number, would have to be divisible by $\pi$. Likewise,
conversely, every real prime $n$ is also divisible by at least one
prime $\pi$. Let, then, $n = \pi\alpha$, from which, by forming
the norm, $n^2 = N(\pi)N(\alpha)$ follows. Since $n$ is a prime and
$N(\pi)N(\alpha)$ integers, two cases are possible:
\[
\begin{aligned}
1.\quad & N(\pi) = n, \quad N(\alpha) = n,\\
2.\quad & N(\pi) = n^2, \quad N(\alpha) = 1.
\end{aligned}
\]

In the first case, if $\pi'$ denotes the number conjugate to $\pi$
and $\pi = a+bi$ is set, then $n = \pi\pi' = a^2+b^2$,
$\alpha = \pi'$, and one sees that this case occurs only when
$n$ belongs to the primes $p$ that are decomposable into the sum of two
squares. Conversely, every such number $p$
can be decomposed into two conjugate factors $\pi\pi'$, of which neither
is a unit. \gesperrt{The primes} $p$ \gesperrt{are therefore not primes in the field}
$R(i)$.

In the second case, which accordingly occurs with the primes $q$,
$\alpha$ is a unit, hence $n$ is associated with $\pi$; that is, the \gesperrt{real
primes} $q$ \gesperrt{are also primes in the field} $R(i)$.

The prime 2 belongs, as already remarked, to the first
kind, and $2 = (1+i)(1-i)$. But it occupies a special
position, because $1-i = -i(1+i)$, and hence $2 = -i(1+i)^2$.
The real prime 2 is therefore (apart from the factor $-i$)
\gesperrt{the square of a prime} in the field $R(i)$. This
case, as one easily sees, occurs with none of the remaining
primes $p$.

% [sic] German authority prints "des Körper R(i)" for grammatical "des Körpers R(i)"; diplomatic DE preserved.
The totality of the primes of the field $R(i)$ thus consists
of the real primes $q$ and of the factors of the primes
$p$. If $p$ is decomposed into the sum of two squares, $a^2+b^2$, then
one also knows the conjugate factors $a\pm bi$ of $p$. To this
%CONT>>
